\chapter{Técnicas avanzadas de demostración y el teorema del emparedado}

\section{La suma de sucesiones convergentes}

Una de las ideas más fundamentales en análisis es que ``las operaciones
agradables preservan la convergencia''. Si dos sucesiones convergen cada
una, entonces su suma también converge, y converge a la suma de sus
límites.

Esto puede parecer obvio a primera vista: si $a(n)$ se acerca a $L$ y
$b(n)$ se acerca a $M$, ¿no debería $a(n)+b(n)$ acercarse a $L+M$? Aunque
la intuición es correcta, hacer esto riguroso requiere cierta habilidad en
el manejo de nuestra definición épsilon-$N$.

La idea clave es la siguiente: si un ingeniero exige que nuestra producción
combinada esté dentro de $\varepsilon$ del objetivo $L+M$, no podemos
simplemente pedir a cada fábrica que cumpla independientemente la
tolerancia completa $\varepsilon$. En cambio, debemos ser inteligentes al
repartir nuestro ``presupuesto de tolerancia''.

Piénselo así: si la primera fábrica puede garantizar que su producción
está dentro de $\varepsilon/2$ de $L$, y la segunda fábrica puede
garantizar que su producción está dentro de $\varepsilon/2$ de $M$,
entonces, por la desigualdad triangular, su suma estará dentro de
$\varepsilon$ de $L+M$. ¡Este es el corazón de la demostración!

\subsection{La configuración matemática}

Supongamos que tenemos:
\begin{itemize}
  \item Una sucesión $a : \mathbb{N} \to \mathbb{R}$ que converge a $L$.
  \item Una sucesión $b : \mathbb{N} \to \mathbb{R}$ que converge a $M$.
  \item Una sucesión $c : \mathbb{N} \to \mathbb{R}$ con la propiedad de
        que $c(n) = a(n) + b(n)$ para todo $n$.
\end{itemize}

Queremos demostrar que $c$ converge a $L+M$.

\subsection{Estrategia}

\begin{enumerate}
  \item Comenzar desplegando las definiciones de convergencia en la meta y
        las hipótesis.
  \item Dado cualquier $\varepsilon > 0$, usar la convergencia de $a$ para
        obtener un $N_a$ que funcione para $\varepsilon/2$.
  \item Similarmente, usar la convergencia de $b$ para obtener un $N_b$
        que funcione para $\varepsilon/2$.
  \item Tomar $N = N_a + N_b$ (asegurando que ambas condiciones se
        satisfagan).
  \item Usar la desigualdad triangular para combinar las dos
        semitolerancias.
\end{enumerate}

\subsection{Solución en Lean}

\begin{minted}{lean}
-- ---------------------------------------------------------------------
-- Demostrar que si la sucesión a converge a L y la b a M, entonces a+b
-- converge a L+M.
-- ----------------------------------------------------------------------

-- Demostración en lenguaje natural
-- ================================

-- Tenemos que probar que para todo ε ∈ ℝ, si
--    ε > 0                                                          (1)
-- entonces
--    (∃N ∈ ℕ)(∀n ∈ ℕ)[n ≥ N → |(a + b)(n) - (L + M)| < ε]           (2)
--
-- Por (1), se tiene que
--    ε/2 > 0                                                        (3)
-- Por (3) y porque el límite de a es L, se tiene que
--    (∃N ∈ ℕ)(∀n ∈ ℕ)[n ≥ N → |a(n) - L| < ε/2]
-- Sea N₁ ∈ ℕ tal que
--    (∀n ∈ ℕ)[n ≥ N₁ → |a(n) - L| < ε/2]                            (4)
-- Por (3) y porque el límite de b es M, se tiene que
--    (∃N ∈ ℕ)(∀n ∈ ℕ)[n ≥ N → |b(n) - M| < ε/2]
-- Sea N₂ ∈ ℕ tal que
--    (∀n ∈ ℕ)[n ≥ N₂ → |b(n) - M| < ε/2]                            (5)
-- Sea N = máx(N₁, N₂). Veamos que verifica la condición (1). Para ello,
-- sea n ∈ ℕ tal que n ≥ N. Entonces, n ≥ N₁ y n ≥ N₂. Por tanto, por
-- las propiedades (4) y (5) se tiene que
--    |a(n) - L| < ε/2                                               (6)
--    |b(n) - M| < ε/2                                               (7)
-- Finalmente,
--    |(a + b)(n) - (L + M)| = |(a(n) + b(n)) - (L + M)|
--                           = |(a(n) - L) + (b(n) - M)|
--                           ≤ |a(n) - L| + |b(n) - M|
--                           < ε / 2 + ε / 2                [por (6) y (7)
--                           = ε

-- Demostraciones con Lean4
-- ========================

import Mathlib.Data.Real.Basic
import Mathlib.Tactic
import Analisis_en_Lean.C2_El_calculo_de_pi

variable {a b c : ℕ → ℝ}
variable {L M : ℝ}

-- 1ª demostración
-- ===============

namespace Solucion1

example
  (ha : LimSuc a L)
  (hb : LimSuc b M)
  : LimSuc (a + b) (L + M) :=
by
  intros ε hε
  -- ε : ℝ
  -- hε : ε > 0
  -- ⊢ ∃ N, ∀ (n : ℕ), n ≥ N → |(a + b) n - (L + M)| < ε
  obtain ⟨Na, hNa⟩ := ha (ε / 2) (by linarith)
  -- Na : ℕ
  -- hNa : ∀ (n : ℕ), n ≥ Na → |a n - L| < ε / 2
  obtain ⟨Nb, hNb⟩ := hb (ε / 2) (by linarith)
  -- Nb : ℕ
  -- hNb : ∀ (n : ℕ), n ≥ Nb → |b n - M| < ε / 2
  use max Na Nb
  -- ⊢ ∀ (n : ℕ), n ≥ max Na Nb → |(s + t) n - (L + M)| < ε
  intros n hn
  -- n : ℕ
  -- hn : n ≥ N
  -- ⊢ |(a + b) n - (L + M)| < ε
  calc |(a + b) n - (L + M)|
     = |(a n + b n) - (L + M)|  := rfl
   _ = |(a n - L) + (b n - M)|  := by ring_nf
   _ ≤ |a n - L| + |b n - M|    := abs_add_le _ _
   _ < ε / 2 + ε / 2            := add_lt_add (hNa n (le_of_max_le_left hn))
                                              (hNb n (le_of_max_le_right hn))
   _ = ε                        := add_halves ε

end Solucion1

-- 2ª demostración
-- ===============

namespace Solucion2

example
  (ha : LimSuc a L)
  (hb : LimSuc b M)
  : LimSuc (a + b) (L + M) :=
by
  intros ε hε
  -- ε : ℝ
  -- hε : ε > 0
  -- ⊢ ∃ N, ∀ (n : ℕ), n ≥ N → |(a + b) n - (L + M)| < ε
  obtain ⟨Na, hNa⟩ := ha (ε / 2) (half_pos hε)
  -- Na : ℕ
  -- hNa : ∀ (n : ℕ), n ≥ Na → |a n - L| < ε / 2
  obtain ⟨Nb, hNb⟩ := hb (ε / 2) (half_pos hε)
  -- Nb : ℕ
  -- hNb : ∀ (n : ℕ), n ≥ Nb → |b n - M| < ε / 2
  refine ⟨max Na Nb, fun n hn => ?_⟩
  -- n : ℕ
  -- hn : n ≥ max Na Nb
  -- ⊢ |(a + b) n - (L + M)| < ε
  have nNa : n ≥ Na := le_of_max_le_left hn
  specialize hNa n nNa
  -- hNa : |a n - L| < ε / 2
  have nNb : n ≥ Nb := le_of_max_le_right hn
  specialize hNb n nNb
  -- hNb : |b n - M| < ε / 2
  calc |(a + b) n - (L + M)|
     = |(a n + b n) - (L + M)|  := rfl
   _ = |(a n - L) + (b n - M)|  := by ring_nf
   _ ≤ |a n - L| + |b n - M|    := by apply abs_add_le
   _ < ε / 2 + ε / 2            := by linarith [hNa, hNb]
   _ = ε                        := by apply add_halves

end Solucion2

-- 3ª demostración
-- ===============

namespace Solucion3

example
  (ha : LimSuc a L)
  (hb : LimSuc b M)
  : LimSuc (a + b) (L + M) :=
by
  intros ε hε
  -- ε : ℝ
  -- hε : ε > 0
  -- ⊢ ∃ N, ∀ n ≥ N, |(a + b) n - (L + M)| < ε
  obtain ⟨Na, hNa⟩ := ha (ε/2) (by linarith)
  -- Na : ℕ
  -- hNa : ∀ n ≥ Na, |a n - L| < ε / 2
  obtain ⟨Nb, hNb⟩ := hb (ε/2) (by linarith)
  -- Nb : ℕ
  -- hNb : ∀ n ≥ Nb, |b n - M| < ε / 2
  use max Na Nb
  -- ⊢ ∀ n ≥ max Na Nb, |(a + b) n - (L + M)| < ε
  intros n hn
  -- n : ℕ
  -- hn : n ≥ max Na Nb
  -- ⊢ |(a + b) n - (L + M)| < ε
  have hn₁ : n ≥ Na := le_of_max_le_left hn
  specialize hNa n hn₁
  -- hNa : |a n - L| < ε / 2
  have hn₂ : n ≥ Nb := le_of_max_le_right hn
  specialize hNb n hn₂
  -- hNb : |b n - M| < ε / 2
  calc |(a + b) n - (L + M)|
       = |(a n + b n) - (L + M)|  := by rfl
     _ = |(a n - L) + (b n -  M)| := by {congr; ring}
     _ ≤ |a n - L| + |b n -  M|   := by apply abs_add_le
     _ < ε / 2 + ε / 2            := by linarith
     _ = ε                        := by apply add_halves

end Solucion3

-- 4ª demostración
-- ===============

namespace Solucion4

example
  (ha : LimSuc a L)
  (hb : LimSuc b M)
  : LimSuc (a + b) (L + M) :=
by
  intros ε hε
  -- ε : ℝ
  -- hε : ε > 0
  -- ⊢ ∃ N, ∀ n ≥ N, |(a + b) n - (L + M)| < ε
  obtain ⟨Na, hNa⟩ := ha (ε/2) (by linarith)
  -- Na : ℕ
  -- hNa : ∀ n ≥ Na, |a n - L| < ε / 2
  obtain ⟨Nb, hNb⟩ := hb (ε/2) (by linarith)
  -- Nb : ℕ
  -- hNb : ∀ n ≥ Nb, |b n - M| < ε / 2
  use max Na Nb
  -- ⊢ ∀ n ≥ max Na Nb, |(a + b) n - (L + M)| < ε
  intros n hn
  -- n : ℕ
  -- hn : n ≥ max Na Nb
  -- ⊢ |(a + b) n - (L + M)| < ε
  obtain ⟨hn₁, hn₂⟩ := max_le_iff.mp hn
  -- hn₁ : n ≥ Na
  -- hn₂ : n ≥ Nb
  calc |(a + b) n - (L + M)|
       = |a n + b n - (L + M)|   := by rfl
     _ = |(a n - L) + (b n - M)| := by { congr; ring }
     _ ≤ |a n - L| + |b n - M|   := by apply abs_add_le
     _ < ε/2 + ε/2               := add_lt_add (hNa n hn₁) (hNb n hn₂)
     _ = ε                       := by simp

end Solucion4

-- 5ª demostración
-- ===============

namespace Solucion5

example
  (ha : LimSuc a L)
  (hb : LimSuc b M)
  : LimSuc (a + b) (L + M) :=
by
  intros ε hε
  -- ε : ℝ
  -- hε : ε > 0
  -- ⊢ ∃ N, ∀ n ≥ N, |(a + b) n - (L + M)| < ε
  obtain ⟨Na, hNa⟩ := ha (ε/2) (by linarith)
  -- Na : ℕ
  -- hNa : ∀ n ≥ Na, |a n - L| < ε / 2
  obtain ⟨Nb, hNb⟩ := hb (ε/2) (by linarith)
  -- Nb : ℕ
  -- hNb : ∀ n ≥ Nb, |b n - M| < ε / 2
  use max Na Nb
  -- ⊢ ∀ n ≥ max Na Nb, |(a + b) n - (L + M)| < ε
  intros n hn
  -- n : ℕ
  -- hn : n ≥ max Na Nb
  -- ⊢ |(a + b) n - (L + M)| < ε
  obtain ⟨hn₁, hn₂⟩ := max_le_iff.mp hn
  -- hn₁ : Na ≤ n
  -- hn₂ : Nb ≤ n
  have cota₁ : |a n - L| < ε/2 := hNa n hn₁
  have cota₂ : |b n - M| < ε/2 := hNb n hn₂
  calc |(a + b) n - (L + M)|
       = |(a n + b n) - (L + M)| := by rfl
     _ = |(a n - L) + (b n - M)| := by { congr; ring }
     _ ≤ |a n - L| + |b n - M|   := by apply abs_add_le
     _ < ε                       := by linarith

end Solucion5

-- 6ª demostración
-- ===============

namespace Solucion6

example
  (ha : LimSuc a L)
  (hb : LimSuc b M)
  : LimSuc (a + b) (L + M) :=
by
  intros ε Hε
  -- ε : ℝ
  -- Hε : ε > 0
  -- ⊢ ∃ N, ∀ n ≥ N, |(a + b) n - (L + M)| < ε
  obtain ⟨Na, hNa⟩ := ha (ε/2) (by linarith)
  -- Na : ℕ
  -- hNa : ∀ n ≥ Na, |a n - L| < ε / 2
  obtain ⟨Nb, hNb⟩ := hb (ε/2) (by linarith)
  -- Nb : ℕ
  -- hNb : ∀ n ≥ Nb, |b n - M| < ε / 2
  use max Na Nb
  -- ⊢ ∀ n ≥ max Na Nb, |(a + b) n - (L + M)| < ε
  have h1 : max Na Nb ≥ Na := le_max_left _ _
  have h2 : max Na Nb ≥ Nb := le_max_right _ _
  intros n Hn
  -- n : ℕ
  -- Hn : n ≥ max Na Nb
  -- ⊢ |(a + b) n - (L + M)| < ε
  have h3 : |a n - L| < ε/2 := hNa n (by linarith)
  have h4 : |b n - M| < ε/2 := hNb n (by linarith)
  calc |(a + b) n - (L + M)|
       = |(a n + b n) - (L + M)|   := by rfl
     _ = |(a n - L) + (b n - M)|   := by {congr; ring }
     _ ≤ |(a n - L)| + |(b n - M)| := by apply abs_add_le
     _ < ε/2 + ε/2                 := by linarith
     _ = ε                         := by ring

end Solucion6

-- 7ª demostración
-- ===============

namespace Solucion7

example
  (ha : LimSuc a L)
  (hb : LimSuc b M)
  (hc : ∀ n, c n = a n + b n)
  : LimSuc c (L + M) :=
by
  change ∀ ε > 0, ∃ N : ℕ, ∀ n ≥ N, |c n - (L + M)| < ε
  -- ⊢ ∀ ε > 0, ∃ N, ∀ n ≥ N, |c n - (L + M)| < ε
  intro ε hε
  -- ε : ℝ
  -- hε : ε > 0
  -- ⊢ ∃ N, ∀ n ≥ N, |c n - (L + M)| < ε
  change ∀ ε₁ > 0, ∃ Na : ℕ, ∀ n ≥ Na, |a n - L| < ε₁ at ha
  -- ha : ∀ ε₁ > 0, ∃ Na, ∀ n ≥ Na, |a n - L| < ε₁
  change ∀ ε₂ > 0, ∃ Nb : ℕ, ∀ n ≥ Nb, |b n - M| < ε₂ at hb
  -- hb : ∀ ε₂ > 0, ∃ Nb, ∀ n ≥ Nb, |b n - M| < ε₂
  specialize ha (ε / 2)
  -- ha : ε / 2 > 0 → ∃ Na, ∀ n ≥ Na, |a n - L| < ε / 2
  specialize hb (ε / 2)
  -- ⊢ ∃ N, ∀ n ≥ N, |c n - (L + M)| < ε
  have hε2 : 0 < ε / 2 := by linarith [hε]
  specialize ha hε2
  -- ha : ∃ Na, ∀ n ≥ Na, |a n - L| < ε / 2
  specialize hb hε2
  -- hb : ∃ Nb, ∀ n ≥ Nb, |b n - M| < ε / 2
  choose Na hNa using ha
  -- Na : ℕ
  -- hNa : ∀ n ≥ Na, |a n - L| < ε / 2
  choose Nb hNb using hb
  -- Nb : ℕ
  -- hNb : ∀ n ≥ Nb, |b n - M| < ε / 2
  use Na + Nb
  -- ⊢ ∀ n ≥ Na + Nb, |c n - (L + M)| < ε
  intro n hn
  -- n : ℕ
  -- hn : n ≥ Na + Nb
  -- ⊢ |c n - (L + M)| < ε
  specialize hc n
  -- hc : c n = a n + b n
  rewrite [hc]
  -- ⊢ |a n + b n - (L + M)| < ε
  have h1 : a n + b n - (L + M) =
    (a n - L) + (b n - M) := by ring_nf
  rewrite [h1]
  -- ⊢ |a n - L + (b n - M)| < ε
  specialize hNa n
  -- hNa : n ≥ Na → |a n - L| < ε / 2
  specialize hNb n
  -- hNb : n ≥ Nb → |b n - M| < ε / 2
  have h2 : Na ≤ n := by linarith
  have h3 : Nb ≤ n := by linarith
  specialize hNa h2
  -- hNa : |a n - L| < ε / 2
  specialize hNb h3
  -- hNb : |b n - M| < ε / 2
  have h4 : |a n - L + (b n - M)| ≤
    |a n - L| + |(b n - M)| := by apply abs_add_le
  linarith [abs_add_le (a n - L) (b n - M)]

end Solucion7

-- 8ª demostración
-- ===============

namespace Solucion8

example
  (ha : LimSuc a L)
  (hb : LimSuc b M)
  : LimSuc (a + b) (L + M) :=
by
  intros ε hε
  -- ε : ℝ
  -- hε : ε > 0
  -- ⊢ ∃ N, ∀ n ≥ N, |(a + b) n - (L + M)| < ε
  obtain ⟨Na, hNa⟩ := ha (ε/2) (by linarith)
  -- Na : ℕ
  -- hNa : ∀ n ≥ Na, |a n - L| < ε / 2
  obtain ⟨Nb, hNb⟩ := hb (ε/2) (by linarith)
  -- Nb : ℕ
  -- hNb : ∀ n ≥ Nb, |b n - M| < ε / 2
  use max Na Nb
  -- ⊢ ∀ n ≥ max Na Nb, |(a + b) n - (L + M)| < ε
  intros n hn
  -- n : ℕ
  -- hn : n ≥ max Na Nb
  -- ⊢ |(a + b) n - (L + M)| < ε
  specialize hNa n (le_of_max_le_left hn)
  -- hNa : |a n - L| < ε / 2
  specialize hNb n (le_of_max_le_right hn)
  -- hNb : |b n - M| < ε / 2
  simp
  -- ⊢ |a n + b n - (L + M)| < ε
  grind

end Solucion8

-- 9ª demostración
-- ===============

namespace Solucion9

example
  (ha : LimSuc a L)
  (hb : LimSuc b M)
  : LimSuc (a + b) (L + M) :=
by
  intros ε hε
  -- ε : ℝ
  -- hε : ε > 0
  -- ⊢ ∃ N, ∀ (n : ℕ), n ≥ N → |(a + b) n - (L + M)| < ε
  obtain ⟨Na, hNa⟩ := ha (ε / 2) (by grind)
  -- Na : ℕ
  -- hNa : ∀ (n : ℕ), n ≥ Na → |a n - L| < ε / 2
  obtain ⟨Nb, hNb⟩ := hb (ε / 2) (by grind)
  -- Nb : ℕ
  -- hNb : ∀ (n : ℕ), n ≥ Nb → |b n - M| < ε / 2
  use max Na Nb
  -- ⊢ ∀ (n : ℕ), n ≥ max Na Nb → |(s + t) n - (L + M)| < ε
  intros n hn
  -- n : ℕ
  -- hn : n ≥ N
  -- ⊢ |(a + b) n - (L + M)| < ε
  calc |(a + b) n - (L + M)|
     = |(a n + b n) - (L + M)|  := rfl
   _ = |(a n - L) + (b n - M)|  := by grind
   _ ≤ |a n - L| + |b n - M|    := by grind
   _ < ε / 2 + ε / 2            := by grind
   _ = ε                        := by grind

end Solucion9

-- Lemas usados
-- ============

variable (x y z z' : ℝ)
#check (abs_add_le x y : |x + y| ≤ |x| + |y|)
#check (add_halves x : x / 2 + x / 2 = x)
#check (add_lt_add : x < y → z < z' → x + z < y + z')
#check (half_pos : x > 0 → x / 2 > 0)
#check (le_max_left x y : x ≤ max x y)
#check (le_max_right x y : y ≤ max x y)
#check (le_of_max_le_left : max x y ≤ z → x ≤ z)
#check (le_of_max_le_right : max x y ≤ z → y ≤ z)
#check (max_le_iff : max x y ≤ z ↔ x ≤ z ∧ y ≤ z)
\end{minted}

\subsection{Demostración en Lenguaje Natural}

\begin{theorem}
Si dos sucesiones de números reales convergen a sus respectivos límites,
entonces la sucesión formada sumando los términos correspondientes también
converge, y su límite es la suma de los límites originales.
\end{theorem}

\begin{proof}
Supongamos que las sucesiones $a(n)$ y $b(n)$ convergen a $L$ y $M$
respectivamente, y queremos demostrar que $c(n) = a(n)+b(n)$ converge a
$L+M$.

Por definición, debemos demostrar que para cualquier tolerancia
$\varepsilon > 0$, podemos encontrar un punto $N$ tal que para todo
$n \geq N$ se tiene $|c(n)-(L+M)| < \varepsilon$.

Como $a(n)$ converge a $L$, podemos encontrar $N_1$ tal que
$|a(n)-L| < \varepsilon/2$ para todo $n \geq N_1$. Como $b(n)$ converge a
$M$, podemos encontrar $N_2$ tal que $|b(n)-M| < \varepsilon/2$ para todo
$n \geq N_2$.

Sea $N = N_1 + N_2$. Entonces, para cualquier $n \geq N$:
\[
  |c(n)-(L+M)| = |(a(n)+b(n))-(L+M)| = |(a(n)-L)+(b(n)-M)|.
\]
Por la desigualdad triangular, esto es a lo sumo:
\[
  |a(n)-L| + |b(n)-M| < \frac{\varepsilon}{2} + \frac{\varepsilon}{2} = \varepsilon.
\]
Por lo tanto, $c(n)$ converge a $L+M$.
\end{proof}

\section{Separar conjunciones: Descomponiendo metas complejas}

Las demostraciones matemáticas a menudo requieren establecer múltiples
hechos relacionados simultáneamente. Cuando la meta consiste en demostrar
múltiples enunciados conectados por ``y'' ($\wedge$), la táctica
\texttt{split\_ands} resulta invaluable.

Piense en \texttt{split\_ands} como una forma de descomponer una
especificación de fabricación compleja en verificaciones de calidad
individuales. En lugar de intentar verificar que un producto cumple tres
estándares diferentes a la vez, podemos abordar cada estándar de forma
separada y sistemática.

\subsection{Nuevas herramientas}

La táctica \texttt{split\_ands} divide las metas ``y'' en piezas
individuales. Si la meta es $h_1 \wedge h_2 \wedge h_3$, llamar a
\texttt{split\_ands} la dividirá en tres metas separadas: primero $h_1$,
luego $h_2$ y finalmente $h_3$.

\subsection{Ejemplo}

\begin{minted}{lean}
import Mathlib.Data.Real.Basic
import Mathlib.Tactic

-- ---------------------------------------------------------------------
-- Demostrar que si `x = 2` y adenás `y = 3`, entonces se tiene que
--    x = 2 ∧ y = 3
-- ---------------------------------------------------------------------

-- 1ª demostración
-- ===============

example
  (x y : ℝ)
  (hx : x = 2)
  (hy : y = 3)
  : x = 2 ∧ y = 3 :=
And.intro hx hy

-- 2ª demostración
-- ===============

example
  (x y : ℝ)
  (hx : x = 2)
  (hy : y = 3)
  : x = 2 ∧ y = 3 :=
⟨hx, hy⟩

-- 3ª demostración
-- ===============

example
  (x y : ℝ)
  (hx : x = 2)
  (hy : y = 3)
  : x = 2 ∧ y = 3 :=
by
  constructor
  · -- ⊢ x = 2
    exact hx
  · -- hy : y = 3
    exact hy

-- 4ª demostración
-- ===============

example
  (x y : ℝ)
  (hx : x = 2)
  (hy : y = 3)
  : x = 2 ∧ y = 3 :=
by
  split_ands
  · -- ⊢ x = 2
    exact hx
  · -- hy : y = 3
    exact hy

-- 5ª demostración
-- ===============

example
  (x y : ℝ)
  (hx : x = 2)
  (hy : y = 3)
  : x = 2 ∧ y = 3 :=
by
  tauto
\end{minted}

La táctica \texttt{split\_ands} puede parecer sencilla, pero es
increíblemente poderosa para organizar demostraciones complejas. Muchos
teoremas matemáticos importantes tienen conclusiones que son
conjunciones---enunciados de la forma ``A y B y C''. Poder
descomponerlos sistemáticamente hace que las demostraciones sean mucho más
manejables y legibles.

\section{Izquierda y derecha: Tomando decisiones en matemáticas}

Tras dominar \texttt{split\_ands} para manejar situaciones en las que
necesitamos demostrar múltiples cosas simultáneamente, nos dirigimos ahora
a un escenario fundamentalmente diferente: demostrar que al menos una de
varias posibilidades es verdadera. Este es el mundo de los enunciados
``o'' ($\vee$).

Mientras que \texttt{split\_ands} consistía en ser exhaustivos---demostrando
cada parte de una conjunción---, demostrar un enunciado ``o'' (disyunción)
consiste en hacer una elección estratégica. ¡Cuando se trata de demostrar
``$P$ o $Q$'', no es necesario demostrar ambos $P$ y $Q$, basta con
demostrar uno de ellos!

\subsection{Nuevas herramientas}

Cuando la meta es demostrar un enunciado ``o'', $P \vee Q$, se puede hacer
demostrando $P$ o $Q$:
\begin{itemize}
  \item Si se quiere demostrar $P$, se escribe \texttt{left} y la meta se
        convierte en $P$.
  \item Si se quiere demostrar $Q$, se escribe \texttt{right} y la meta se
        convierte en $Q$.
\end{itemize}

\subsection{Ejemplo}

\begin{minted}{lean}
import Mathlib.Data.Real.Basic
import Mathlib.Tactic

-- ---------------------------------------------------------------------
-- Demostrar que si `y = 3`, entonces se tiene que
--    x = 3 ∨ y = 3
-- ---------------------------------------------------------------------

-- 1ª demostración
-- ===============

example
  (x y : ℝ)
  (hy : y = 3)
  : x = 3 ∨ y = 3 :=
Or.inr hy

-- 2ª demostración
-- ===============

example
  (x y : ℝ)
  (hy : y = 3)
  : x = 3 ∨ y = 3 :=
by
  right
  -- ⊢ y = 3
  exact hy

-- 3ª demostración
-- ===============

example
  (x y : ℝ)
  (hy : y = 3)
  : x = 3 ∨ y = 3 :=
by
  tauto
\end{minted}

Los enunciados ``o'' están en todas partes en matemáticas. Las
demostraciones de existencia a menudo tienen este sabor, y poder elegir la
rama correcta puede simplificar drásticamente la estrategia de demostración.

\section{Notación de punto: Accediendo a partes de información compleja}

A menudo en matemáticas se nos dará una hipótesis que contiene múltiples
piezas de información agrupadas. Por ejemplo, podríamos saber que ``$x=2$
Y $y=3$'' pero solo necesitar el hecho de que ``$y=3$'' para nuestra meta
actual. Lean proporciona una forma elegante: la notación de punto.

\subsection{Reglas de la notación de punto}

Cuando se tiene una hipótesis $h : P \wedge Q$, se puede acceder a:
\begin{itemize}
  \item La primera parte con \texttt{h.1} (que da $P$).
  \item La segunda parte con \texttt{h.2} (que da $Q$).
\end{itemize}

Para conjunciones más largas como $P \wedge Q \wedge R$, nótese que
\texttt{h.1} da $P$, pero \texttt{h.2} da $Q \wedge R$ debido a paréntesis
ocultos: $P \wedge Q \wedge R = P \wedge (Q \wedge R)$. Para obtener $Q$
solo, se usa \texttt{h.2.1}, y para $R$ se usa \texttt{h.2.2}.

\subsection{Ejemplo}

\begin{minted}[linenos, breaklines, fontsize=\small]{lean}
import Mathlib.Data.Real.Basic
import Mathlib.Tactic

-- ---------------------------------------------------------------------
-- Demostrar que si `x = 2 ∧ y = 3`, entonces `y = 3`.
-- ---------------------------------------------------------------------

-- 1ª demostración
-- ===============

example
  (x y : ℝ)
  (h : x = 2 ∧ y = 3)
  : y = 3 :=
h.right

-- 2ª demostración
-- ===============

example
  (x y : ℝ)
  (h : x = 2 ∧ y = 3)
  : y = 3 :=
h.2

-- 3ª demostración
-- ===============

example
  (x y : ℝ)
  (h : x = 2 ∧ y = 3)
  : y = 3 :=
by
  exact h.2

-- 4ª demostración
-- ===============

example
  (x y : ℝ)
  (h : x = 2 ∧ y = 3)
  : y = 3 :=
by
  tauto
\end{minted}

La notación de punto es esencial para mantener la claridad en demostraciones
complejas donde múltiples condiciones o propiedades están agrupadas.

\section{\texttt{cases'}: Manejando todas las posibilidades}

Cuando se tiene una hipótesis $h : P \vee Q$, se sabe que $P$ es verdadera
o $Q$ es verdadera, pero no cuál de ellas. Para continuar con la
demostración, es necesario considerar ambas posibilidades de forma
sistemática. La táctica \texttt{cases'} hace exactamente esto.

\subsection{La Táctica \texttt{cases'}}

Cuando se tiene una hipótesis $h : P \vee Q$, se puede escribir
\texttt{cases' h with h1 h2}. Esto crea dos metas separadas:
\begin{itemize}
  \item En la primera meta, se obtiene una nueva hipótesis $h_1 : P$.
  \item En la segunda meta, se obtiene una nueva hipótesis $h_2 : Q$.
\end{itemize}
Deben resolverse ambas metas para completar la demostración, asegurando
que se han cubierto todas las posibilidades lógicas.

\subsection{Ejemplo}

\begin{minted}{lean}
import Mathlib.Data.Real.Basic
import Mathlib.Tactic

-- ---------------------------------------------------------------------
-- Demostrar que si `x = 2 ∨ y = 3`, entonces `(x - 2)(y - 3) = 0`.
-- ---------------------------------------------------------------------

-- 1ª demostración
-- ===============

example
  (x y : ℝ)
  (h : x = 2 ∨ y = 3)
  : (x - 2) * (y - 3) = 0 :=
by
  obtain (h1 | h2) := h
  · -- h1 : x = 2
    rewrite [h1]
    -- ⊢ (2 - 2) * (y - 3) = 0
    simp
  · -- h2 : y = 3
    rewrite [h2]
    -- ⊢ (x - 2) * (3 - 3) = 0
    simp

-- 2ª demostración
-- ===============

example
  (x y : ℝ)
  (h : x = 2 ∨ y = 3)
  : (x - 2) * (y - 3) = 0 :=
by
  rcases h with (h1 | h2)
  · -- h1 : x = 2
    rewrite [h1]
    -- ⊢ (2 - 2) * (y - 3) = 0
    simp
  · -- h2 : y = 3
    rewrite [h2]
    -- ⊢ (x - 2) * (3 - 3) = 0
    simp

-- 3ª demostración
-- ===============

example
  (x y : ℝ)
  (h : x = 2 ∨ y = 3)
  : (x - 2) * (y - 3) = 0 :=
by
  cases h with
  | inl h1 =>
    -- h1 : x = 2
    rewrite [h1]
    -- ⊢ (2 - 2) * (y - 3) = 0
    simp
  | inr h2 =>
    -- h2 : y = 3
    rewrite [h2]
    -- ⊢ (x - 2) * (3 - 3) = 0
    simp

-- 4ª demostración
-- ===============

example
  (x y : ℝ)
  (h : x = 2 ∨ y = 3)
  : (x - 2) * (y - 3) = 0 :=
by
  grind
\end{minted}

\subsection{El conjunto completo de herramientas y/o}

Ahora se dispone de la matriz completa de Y/O:

\begin{center}
\begin{tabular}{|l|c|c|}
\hline
 & $\wedge$ & $\vee$ \\
\hline
\textbf{Meta} & \texttt{split\_ands} & \texttt{left}/\texttt{right} \\
\textbf{Hipótesis} & \texttt{h.1}, \texttt{h.2} & \texttt{cases'} \\
\hline
\end{tabular}
\end{center}

El análisis por casos está en todas partes en matemáticas, desde demostrar
que todo entero es par o impar, hasta mostrar que las funciones continuas
en intervalos cerrados alcanzan sus extremos.

\section{\texttt{abs\_lt}: Trabajando con valores absolutos en convergencia}

Trabajar con valores absolutos es fundamental en análisis real,
especialmente en el contexto de la convergencia de sucesiones. A veces
necesitamos extraer información más específica de las condiciones de valor
absoluto, como cotas direccionales.

\subsection{El teorema \texttt{abs\_lt}}

El teorema \texttt{abs\_lt} establece que $|x| < y$ si y solo si
$-y < x \wedge x < y$. Esto permite convertir entre desigualdades de valor
absoluto y conjunciones de desigualdades regulares, haciéndolas más fáciles
de trabajar en las demostraciones.

La idea clave es que $|x| < y$ captura la idea de que $x$ está dentro de
la distancia $y$ de cero, lo que significa que $x$ pertenece al intervalo
$(-y, y)$.

\subsection{Ejemplo de aplicación}

\begin{minted}{lean}
-- ---------------------------------------------------------------------
-- Demostrar que si la sucesión a es convergente, entonces a está
-- eventualmente acotadas inferiormente; es decir, si a converge a L,
-- entonces
--    (∃ N)(∀ n ≥ N)[aₙ ≥ L - 1]
-- ---------------------------------------------------------------------

-- Demostración en lenguaje natural
-- ================================

-- Puesto que a converge a L, existe un N ∈ ℕ tal que
--    (∀ n ≥ N) |aₙ - L| < 1                                          (1)
-- Veamos que
--    (∀ n ≥ N) [aₙ ≥ L - 1]
-- En efecto, sea n ≥ N. Por (1),
--    |aₙ - L| < 1
-- Luego,
--    -1 < aₙ - L ∧ aₙ - L < 1
--    ==> -1 < aₙ - L
--    ==> aₙ ≥ L - 1

-- Demostraciones con Lean4
-- ========================

import Mathlib.Data.Real.Basic
import Mathlib.Tactic
import Analisis_en_Lean.C2_El_calculo_de_pi

variable {a : ℕ → ℝ}
variable {L : ℝ}

-- 1ª demostración
-- ===============

namespace Solucion1

variable (x : ℝ)

lemma L1
  (h : -1 < x)
  : 0 ≤ x + 1 :=
by
  have h1 : -1 + 1 < x + 1 := (add_lt_add_iff_right 1).mpr h
  have h2 : 0 < x + 1 := neg_lt_iff_pos_add.mp h
  exact le_of_lt h2

example
  (ha : LimSuc a L) :
  ∃ N, ∀ n ≥ N, a n ≥ L - 1 :=
by
  obtain ⟨N, hN⟩ := ha 1 one_pos
  -- N : ℕ
  -- hN : ∀ n ≥ N, |a n - L| < 1
  use N
  -- ⊢ ∀ n ≥ N, a n ≥ L - 1
  intro n hn
  -- n : ℕ
  -- hn : n ≥ N
  -- ⊢ a n ≥ L - 1
  have h1 : |a n - L| < 1 := hN n hn
  have h2 : -1 < a n - L ∧ a n - L < 1 := abs_lt.mp h1
  have h3 : -1 < a n - L := h2.1
  rw [ge_iff_le, ← sub_nonneg]
  -- ⊢ 0 ≤ a n - (L - 1)
  calc 0
     _ ≤ (a n - L) + 1 := L1 (a n - L) h3
     _ = a n - (L - 1) := sub_add (a n) L 1

end Solucion1

-- 2ª demostración
-- ===============

namespace Solucion2

example
  (ha : LimSuc a L) :
  ∃ N, ∀ n ≥ N, a n ≥ L - 1 :=
by
  obtain ⟨N, hN⟩ := ha 1 one_pos
  -- N : ℕ
  -- hN : ∀ n ≥ N, |a n - L| < 1
  use N
  -- ⊢ ∀ n ≥ N, a n ≥ L - 1
  intro n hn
  -- n : ℕ
  -- hn : n ≥ N
  -- ⊢ a n ≥ L - 1
  have h1 : -1 < a n - L := (abs_lt.mp (hN n hn)).1
  have h2 : a n - (L - 1) = (a n - L) + 1 := (sub_add (a n) L 1).symm
  rw [ge_iff_le, ← sub_nonneg]
  -- ⊢ 0 ≤ a n - (L - 1)
  rw [h2]
  -- ⊢ 0 ≤ a n - L + 1
  exact neg_le_iff_add_nonneg.mp (le_of_lt h1)

end Solucion2

-- 3ª demostración
-- ===============

namespace Solucion3

example
  (ha : LimSuc a L) :
  ∃ N, ∀ n ≥ N, a n ≥ L - 1 :=
by
  obtain ⟨N, hN⟩ := ha 1 one_pos
  -- N : ℕ
  -- hN : ∀ n ≥ N, |a n - L| < 1
  use N
  -- ⊢ ∀ n ≥ N, a n ≥ L - 1
  intro n hn
  -- n : ℕ
  -- hn : n ≥ N
  -- ⊢ a n ≥ L - 1
  have h1 : -1 < a n - L := (abs_lt.mp (hN n hn)).1
  rw [ge_iff_le, ← sub_nonneg]
  -- ⊢ 0 ≤ a n - (L - 1)
  calc 0
     _ ≤ (a n - L) + 1 := by grind
     _ = a n - (L - 1) := by grind

end Solucion3

-- 4ª demostración
-- ===============

namespace Solucion4

example
  (ha : LimSuc a L) :
  ∃ N, ∀ n ≥ N, a n ≥ L - 1 :=
by
  obtain ⟨N, hN⟩ := ha 1 one_pos
  -- N : ℕ
  -- hN : ∀ n ≥ N, |a n - L| < 1
  use N
  -- ⊢ ∀ n ≥ N, a n ≥ L - 1
  intro n hn
  -- n : ℕ
  -- hn : n ≥ N
  -- ⊢ a n ≥ L - 1
  have : -1 < a n - L := (abs_lt.mp (hN n hn)).1
  linarith

end Solucion4

-- 5ª demostración
-- ===============

namespace Solucion5

example
  (ha : LimSuc a L) :
  ∃ N, ∀ n ≥ N, a n ≥ L - 1 :=
by
  obtain ⟨N, hN⟩ := ha 1 (by grind)
  -- N : ℕ
  -- hN : ∀ n ≥ N, |a n - L| < 1
  refine ⟨N, fun n hn => ?_⟩
  -- n : ℕ
  -- hn : n ≥ N
  -- ⊢ a n ≥ L - 1
  grind

end Solucion5

-- 6ª demostración
-- ===============

namespace Solucion6

example
  (ha : LimSuc a L) :
  ∃ N, ∀ n ≥ N, a n ≥ L - 1 :=
let ⟨N, hN⟩ := ha 1 (by grind)
⟨N, fun n hn => by grind⟩

end Solucion6
\end{minted}

Esta demostración ilustra la extracción de una cota inferior a partir de
una condición de convergencia. Mostramos que cualquier sucesión convergente
está eventualmente acotada inferiormente (relativa a su límite), lo cual es
un bloque de construcción para muchos teoremas importantes.

\section{El teorema del emparedado}

El teorema del emparedado (también conocido como teorema del sándwich o
teorema de compresión) captura de forma hermosa la idea intuitiva de que si
se atrapa una sucesión entre otras dos sucesiones que convergen al mismo
límite, entonces la sucesión atrapada también debe converger a ese límite.

\subsection{La imagen intuitiva}

Imagine tres corredores en una pista. El corredor A y el corredor C se
dirigen ambos a la misma meta $L$, y el corredor B está siempre entre
ellos. No importa cómo A y C zigzagueen de un lado a otro: mientras ambos
lleguen a $L$ y B permanezca entre ellos, B también debe llegar a $L$.
¡Simplemente no hay otro lugar al que ir!

\subsection{El enunciado matemático}

\begin{theorem}[Teorema del emparedado]
Si $a, c : \mathbb{N} \to \mathbb{R}$ convergen ambas a $L$, y $b$ es otra
sucesión comprimida entre $a$ y $c$ (es decir, $a(n) \leq b(n) \leq c(n)$
para todo $n$), entonces $b$ también converge a $L$.
\end{theorem}

\subsection{Demostración en Lean}

\begin{minted}{lean}
-- ---------------------------------------------------------------------
-- Demostrar que si las sucesiones a y c convergen a L, y b está
-- comprimida entre a y c (es decir, aₙ ≤ bₙ ≤ cₙ para todo n), entonces
-- b también converge a L.
-- ---------------------------------------------------------------------

-- Demostración en lenguaje natural
-- ================================

-- Sea ε > 0. Tenemos que demostrar que existe un N tal que,
--    ∀ n ≥ N, |bₙ - L| < ε                                          (1)
--
-- Puesto que a y c convergen a L, existen N₁ y ℕ₂ tales que
--    ∀ n ≥ N₁, |aₙ - L| < ε                                         (2)
--    ∀ n ≥ N₂, |cₙ - L| < ε                                         (3)
-- Sea
--    N = máx(N₁, ℕ₂).                                               (4)
-- Para demostrar (1), sea n ≥ N. Por  (4), se tiene que
--    n ≥ N₁
--    n ≥ N₂
-- Luego, usando (2) y (3), se tiene que
--    |aₙ - L| < ε
--    |cₙ - L| < ε
-- de donde se deduce que
--    - ε < aₙ - L < ε                                               (5)
--    - ε < bₙ - L < ε                                               (6)
-- Luego,
--    -ε < aₙ - L   [por (5)]
--       ≤ bₙ - L   [por hipótesis de b]
--       ≤ cₙ - L   [por hipótesis de b]
--       < ε        [por (6)]
-- Por tanto,
--    -ε < bₙ - L < ε
-- y, finalmente,
--    |bₙ - L| < ε

-- Demostraciones con Lean4
-- ========================

import Mathlib.Data.Real.Basic
import Mathlib.Tactic

variable {a b c : ℕ → ℝ}
variable {L : ℝ}

def LimSuc (a : ℕ → ℝ) (L : ℝ) : Prop :=
  ∀ ε > 0, ∃ N : ℕ, ∀ n ≥ N, |a n - L| < ε

-- 1ª demostración
-- ===============

example
  (ha : LimSuc a L)
  (hc : LimSuc c L)
  (hab : ∀ n, a n ≤ b n)
  (hbc : ∀ n, b n ≤ c n)
  : LimSuc b L :=
by
  intro ε hε
  -- ε : ℝ
  -- hε : ε > 0
  -- ⊢ ∃ N, ∀ n ≥ N, |b n - L| < ε
  obtain ⟨Na, hNa⟩ := ha ε hε
  -- Na : ℕ
  -- hNa : ∀ n ≥ Na, |a n - L| < ε
  obtain ⟨Nc, hNc⟩ := hc ε hε
  -- Nc : ℕ
  -- hNc : ∀ n ≥ Nc, |c n - L| < ε
  exact ⟨max Na Nc, by grind⟩

-- 2ª demostración
-- ===============

example
  (ha : LimSuc a L)
  (hc : LimSuc c L)
  (hab : ∀ n, a n ≤ b n)
  (hbc : ∀ n, b n ≤ c n)
  : LimSuc b L :=
by
  intro ε hε
  -- ε : ℝ
  -- hε : ε > 0
  -- ⊢ ∃ N, ∀ n ≥ N, |b n - L| < ε
  have ha : ∃ N, ∀ n ≥ N, |a n - L| < ε := ha ε hε
  have hb : ∃ N, ∀ n ≥ N, |c n - L| < ε := hc ε hε
  obtain ⟨Na, hNa⟩ := ha
  -- Na : ℕ
  -- hNa : ∀ n ≥ Na, |a n - L| < ε
  obtain ⟨Nc, hNc⟩ := hb
  -- Nc : ℕ
  -- hNc : ∀ n ≥ Nc, |c n - L| < ε
  use max Na Nc
  -- ⊢ ∀ n ≥ max Na Nc, |b n - L| < ε
  intro n hn
  -- n : ℕ
  -- hn : n ≥ max Na Nc
  -- ⊢ |b n - L| < ε
  rewrite [abs_lt]
  -- ⊢ -ε < b n - L ∧ b n - L < ε
  constructor
  · -- ⊢ -ε < b n - L
    calc -ε
       < a n - L := by grind
     _ ≤ b n - L := by grind
  · -- ⊢ b n - L < ε
    calc b n - L
       ≤ c n - L := by grind
     _ < ε       := by grind

-- 3ª demostración
-- ===============

example
    (ha : LimSuc a L)
    (hc : LimSuc c L)
    (hab : ∀ n, a n ≤ b n)
    (hbc : ∀ n, b n ≤ c n) :
    LimSuc b L :=
by
  intro ε hε
  -- ε : ℝ
  -- hε : ε > 0
  -- ⊢ ∃ N, ∀ n ≥ N, |b n - L| < ε
  obtain ⟨Na, hNa⟩ := ha ε hε
  -- Na : ℕ
  -- hNa : ∀ n ≥ Na, |a n - L| < ε
  obtain ⟨Nc, hNc⟩ := hc ε hε
  -- Nc : ℕ
  -- hNc : ∀ n ≥ Nc, |c n - L| < ε
  refine ⟨max Na Nc, fun n hn => ?_⟩
  -- n : ℕ
  -- hn : n ≥ max Na Nc
  -- ⊢ |b n - L| < ε
  have hna := hNa n (le_of_max_le_left hn)
  have hnc := hNc n (le_of_max_le_right hn)
  rw [abs_lt] at hna hnc ⊢
  -- hna : -ε < a n - L ∧ a n - L < ε
  -- hnc : -ε < c n - L ∧ c n - L < ε
  -- ⊢ -ε < b n - L ∧ b n - L < ε
  exact ⟨by linarith [hab n], by linarith [hbc n]⟩

-- 4ª demostración
-- ===============

example
    (ha : LimSuc a L)
    (hc : LimSuc c L)
    (hab : ∀ n, a n ≤ b n)
    (hbc : ∀ n, b n ≤ c n)
    : LimSuc b L :=
by
  intro ε hε
  -- ε : ℝ
  -- hε : ε > 0
  -- ⊢ ∃ N, ∀ n ≥ N, |b n - L| < ε
  obtain ⟨Na, hNa⟩ := ha ε hε
  -- Na : ℕ
  -- hNa : ∀ n ≥ Na, |a n - L| < ε
  obtain ⟨Nc, hNc⟩ := hc ε hε
  -- Nc : ℕ
  -- hNc : ∀ n ≥ Nc, |c n - L| < ε
  refine ⟨max Na Nc, fun n hn => ?_⟩
  -- n : ℕ
  -- hn : n ≥ max Na Nc
  -- ⊢ |b n - L| < ε
  have hna := hNa n (le_of_max_le_left hn)
  have hnc := hNc n (le_of_max_le_right hn)
  rw [abs_lt] at hna hnc ⊢
  -- hna : -ε < a n - L ∧ a n - L < ε
  -- hnc : -ε < c n - L ∧ c n - L < ε
  -- ⊢ -ε < b n - L ∧ b n - L < ε
  refine ⟨?_, ?_⟩
  · -- ⊢ -ε < b n - L
    calc -ε
         < a n - L := hna.1
       _ ≤ b n - L := by linarith [hab n]
  · calc b n - L
         ≤ c n - L := by linarith [hbc n]
       _ < ε       := hnc.2

-- 5ª demostración
-- ===============

example
  (ha : LimSuc a L)
  (hc : LimSuc c L)
  (hab : ∀ n, a n ≤ b n)
  (hbc : ∀ n, b n ≤ c n)
  : LimSuc b L :=
by
  intro ε hε
  -- ε : ℝ
  -- hε : ε > 0
  -- ⊢ ∃ N, ∀ n ≥ N, |b n - L| < ε
  have ha : ∃ N, ∀ n ≥ N, |a n - L| < ε := ha ε hε
  have hb : ∃ N, ∀ n ≥ N, |c n - L| < ε := hc ε hε
  obtain ⟨Na, hNa⟩ := ha
  -- Na : ℕ
  -- hNa : ∀ n ≥ Na, |a n - L| < ε
  obtain ⟨Nc, hNc⟩ := hb
  -- Nc : ℕ
  -- hNc : ∀ n ≥ Nc, |c n - L| < ε
  use max Na Nc
  -- ⊢ ∀ n ≥ max Na Nc, |b n - L| < ε
  intro n hn
  -- n : ℕ
  -- hn : n ≥ max Na Nc
  -- ⊢ |b n - L| < ε
  have hna : Na ≤ n := le_of_max_le_left hn
  have hnc : Nc ≤ n := le_of_max_le_right hn
  have hNa : |a n - L| < ε := hNa n hna
  have hNc : |c n - L| < ε := hNc n hnc
  rewrite [abs_lt] at hNa
  -- hNa : -ε < a n - L ∧ a n - L < ε
  rewrite [abs_lt] at hNc
  -- hNc : -ε < c n - L ∧ c n - L < ε
  rewrite [abs_lt]
  -- ⊢ -ε < b n - L ∧ b n - L < ε
  constructor
  · -- ⊢ -ε < b n - L
    calc -ε
       < a n - L := hNa.1
     _ ≤ b n - L := by simp only [tsub_le_iff_right, sub_add_cancel, hab]
  · -- ⊢ b n - L < ε
    calc b n - L
       ≤ c n - L := by simp only [tsub_le_iff_right, sub_add_cancel, hbc]
     _ < ε       := hNc.2

-- 6ª demostración
-- ===============

example
  (ha : LimSuc a L)
  (hc : LimSuc c L)
  (hab : ∀ n, a n ≤ b n)
  (hbc : ∀ n, b n ≤ c n)
  : LimSuc b L :=
by
  intro ε hε
  -- ε : ℝ
  -- hε : ε > 0
  -- ⊢ ∃ N, ∀ n ≥ N, |b n - L| < ε
  have ha : ∃ N, ∀ n ≥ N, |a n - L| < ε := ha ε hε
  have hb : ∃ N, ∀ n ≥ N, |c n - L| < ε := hc ε hε
  obtain ⟨Na, hNa⟩ := ha
  -- Na : ℕ
  -- hNa : ∀ n ≥ Na, |a n - L| < ε
  obtain ⟨Nc, hNc⟩ := hb
  -- Nc : ℕ
  -- hNc : ∀ n ≥ Nc, |c n - L| < ε
  use max Na Nc
  -- ⊢ ∀ n ≥ max Na Nc, |b n - L| < ε
  intro n hn
  -- n : ℕ
  -- hn : n ≥ max Na Nc
  -- ⊢ |b n - L| < ε
  have hna : Na ≤ n := le_of_max_le_left hn
  have hnc : Nc ≤ n := le_of_max_le_right hn
  have hNa : |a n - L| < ε := hNa n hna
  have hNc : |c n - L| < ε := hNc n hnc
  rewrite [abs_lt] at hNa
  -- hNa : -ε < a n - L ∧ a n - L < ε
  rewrite [abs_lt] at hNc
  -- hNc : -ε < c n - L ∧ c n - L < ε
  rewrite [abs_lt]
  -- ⊢ -ε < b n - L ∧ b n - L < ε
  split_ands
  · -- ⊢ -ε < b n - L
    have hab : a n ≤ b n := hab n
    linarith
  · -- ⊢ b n - L < ε
    have hbc : b n ≤ c n := hbc n
    linarith

-- 7ª demostración
-- ===============

example
    (ha : LimSuc a L)
    (hc : LimSuc c L)
    (hab : ∀ n, a n ≤ b n)
    (hbc : ∀ n, b n ≤ c n)
    : LimSuc b L :=
by
  intro ε hε
  -- ε : ℝ
  -- hε : ε > 0
  -- ⊢ ∃ N, ∀ n ≥ N, |b n - L| < ε
  obtain ⟨Na, hNa⟩ := ha ε hε
  -- Na : ℕ
  -- hNa : ∀ n ≥ Na, |a n - L| < ε
  obtain ⟨Nc, hNc⟩ := hc ε hε
  -- Nc : ℕ
  -- hNc : ∀ n ≥ Nc, |c n - L| < ε
  refine ⟨max Na Nc, ?_⟩
  -- ⊢ ∀ n ≥ max Na Nc, |b n - L| < ε
  intro n hn
  -- n : ℕ
  -- hn : n ≥ max Na Nc
  -- ⊢ |b n - L| < ε
  have hna : n ≥ Na := le_of_max_le_left hn
  have hnc : n ≥ Nc := le_of_max_le_right hn
  have ha' : |a n - L| < ε := hNa n hna
  have hc' : |c n - L| < ε := hNc n hnc
  rw [abs_lt] at ha' hc' ⊢
  -- ⊢ -ε < b n - L ∧ b n - L < ε
  constructor
  · -- ⊢ -ε < b n - L
    have h1 : -ε < a n - L := ha'.1
    have h2 : a n - L ≤ b n - L := by
      linarith [hab n]
    linarith
  · -- ⊢ b n - L < ε
    have h1 : b n - L ≤ c n - L := by
      linarith [hbc n]
    have h2 : c n - L < ε := hc'.2
    linarith

-- Lemas usados
-- ============

variable (a b c : ℝ)
#check (abs_lt : |a| < b ↔ -b < a ∧ a < b)
#check (le_of_max_le_left : max a b ≤ c → a ≤ c)
#check (le_of_max_le_right : max a b ≤ c → b ≤ c)
#check (sub_add_cancel a b : (a - b) + b = a)
#check (tsub_le_iff_right : a - b ≤ c ↔ a ≤ c + b)
\end{minted}

\subsection{Demostración en lenguaje natural}

\begin{proof}
Dado cualquier $\varepsilon > 0$, necesitamos demostrar que
$|b(n)-L| < \varepsilon$ para $n$ suficientemente grande.

Como $a(n) \to L$, existe $N_a$ tal que $|a(n)-L| < \varepsilon$ para todo
$n \geq N_a$. Como $c(n) \to L$, existe $N_c$ tal que $|c(n)-L| < \varepsilon$
para todo $n \geq N_c$.

Sea $N = N_a + N_c$. Para cualquier $n \geq N$, tenemos:
\begin{equation}
  L - \varepsilon < a(n) \leq b(n) \leq c(n) < L + \varepsilon.
\end{equation}
Por lo tanto, $L - \varepsilon < b(n) < L + \varepsilon$, lo que significa
$|b(n)-L| < \varepsilon$. Luego $b(n) \to L$.
\end{proof}

\subsection{Aplicaciones}

El teorema del emparedado es un caballo de batalla del análisis matemático,
utilizado para demostrar resultados de convergencia difíciles reduciéndolos
a problemas más sencillos. Algunos ejemplos son:
\begin{itemize}
  \item $\sin(1/n) \to 0$ (comprimido entre $-1/n$ y $1/n$).
  \item Sucesiones recursivas donde las fórmulas exactas son intratables.
  \item Sucesiones definidas por procesos geométricos o probabilísticos
        complejos.
\end{itemize}

El teorema demuestra el poder de combinar múltiples técnicas: argumentos
épsilon-$N$, manipulación de valores absolutos, descomposición lógica y
razonamiento con desigualdades. Esta síntesis de herramientas es lo que
hace posible la demostración matemática avanzada.

\section{Ejercicios}

\begin{ejercicio}
  Demostrar que si $x=2$, $y=3$ y $z=4$, entonces
  $$x = 2 ∧ y = 3 ∧ z = 4$$
\end{ejercicio}

\begin{minted}{lean}
-- ---------------------------------------------------------------------
-- Demostrar que si `x=2`, `y=3` y `z=4`, entonces
--    x = 2 ∧ y = 3 ∧ z = 4
-- ---------------------------------------------------------------------

import Mathlib.Data.Real.Basic
import Mathlib.Tactic

variable (x y z : ℝ)

-- 1ª demostración
-- ===============

example
  (hx : x = 2)
  (hy : y = 3)
  (hz : z = 4)
  : x = 2 ∧ y = 3 ∧ z = 4 :=
by
  grind

-- 2ª demostración
-- ===============

example
  (hx : x = 2)
  (hy : y = 3)
  (hz : z = 4)
  : x = 2 ∧ y = 3 ∧ z = 4 :=
by
  solve_by_elim

-- 3ª demostración
-- ===============

example
  (hx : x = 2)
  (hy : y = 3)
  (hz : z = 4)
  : x = 2 ∧ y = 3 ∧ z = 4 :=
by
  simp_all

-- 4ª demostración
-- ===============

example
  (hx : x = 2)
  (hy : y = 3)
  (hz : z = 4)
  : x = 2 ∧ y = 3 ∧ z = 4 :=
by
  trivial

-- 5ª demostración
-- ===============

example
  (hx : x = 2)
  (hy : y = 3)
  (hz : z = 4)
  : x = 2 ∧ y = 3 ∧ z = 4 :=
by
  simp [*]

example
  (hx : x = 2)
  (hy : y = 3)
  (hz : z = 4)
  : x = 2 ∧ y = 3 ∧ z = 4 :=
by
  simp only [and_self, hx, hy, hz]

-- 6ª demostración
-- ===============

example
  (hx : x = 2)
  (hy : y = 3)
  (hz : z = 4)
  : x = 2 ∧ y = 3 ∧ z = 4 :=
by
  constructor
  · -- ⊢ x = 2
    exact hx
  · -- ⊢ y = 3 ∧ z = 4
    constructor
    · -- ⊢ y = 3
      exact hy
    · -- ⊢ z = 4
      exact hz

-- 7ª demostración
-- ===============

example
  (hx : x = 2)
  (hy : y = 3)
  (hz : z = 4)
  : x = 2 ∧ y = 3 ∧ z = 4 :=
by
  split_ands
  · -- ⊢ x = 2
    exact hx
  · -- ⊢ y = 3
    exact hy
  · -- ⊢ z = 4
    exact hz

-- 8ª demostración
-- ===============

example
  (hx : x = 2)
  (hy : y = 3)
  (hz : z = 4)
  : x = 2 ∧ y = 3 ∧ z = 4 :=
And.intro hx (And.intro hy hz)

-- 9ª demostración
-- ===============

example
  (hx : x = 2)
  (hy : y = 3)
  (hz : z = 4)
  : x = 2 ∧ y = 3 ∧ z = 4 :=
⟨hx, ⟨hy, hz⟩⟩
\end{minted}

\begin{ejercicio}
  Demostrar que si $x = 2 ∧ y = 3 ∧ z = 4$, entonces $z=4$.
\end{ejercicio}


\begin{minted}{lean}
import Mathlib.Data.Real.Basic
import Mathlib.Tactic

variable (x y z : ℝ)

-- 1ª demostración
-- ===============

example
  (h : x = 2 ∧ y = 3 ∧ z = 4)
  : z = 4 :=
by
  grind

-- 2ª demostración
-- ===============

example
  (h : x = 2 ∧ y = 3 ∧ z = 4)
  : z = 4 :=
by
  simp_all

-- 3ª demostración
-- ===============

example
  (h : x = 2 ∧ y = 3 ∧ z = 4)
  : z = 4 :=
by
  simp [*]

-- 4ª demostración
-- ===============

example
  (h : x = 2 ∧ y = 3 ∧ z = 4)
  : z = 4 :=
by
  simp only [h]

-- 4ª demostración
-- ===============

example
  (h : x = 2 ∧ y = 3 ∧ z = 4)
  : z = 4 :=
h.2.2
\end{minted}

\begin{ejercicio}
  Demostrar que si la sucesión $a$ converge a 5, entonces es
  eventualmente positiva; es decir, existe un $N$ tal que para todo $n ≥ N$
  se tiene que $aₙ > 0$.
\end{ejercicio}

\begin{minted}{lean}
import Mathlib.Data.Real.Basic
import Mathlib.Tactic

-- Demostración en lenguaje natural
-- ================================

-- Puesto que a converge a 5, existe un N ∈ ℕ tal que
--    (∀ n ≥ N) |aₙ - 5| < 5                                          (1)
-- Veamos que
--    (∀ n ≥ N) [aₙ > 0]
-- En efecto, sea n ≥ N. Por (1),
--    |aₙ - 5| < 5
-- Luego,
--    -5 < aₙ - 5 ∧ aₙ - 5 < 5
--    ==> -5 < aₙ - 5
--    ==> aₙ > 0

-- Demostraciones con Lean4
-- ========================

import Mathlib.Data.Real.Basic
import Mathlib.Tactic
import Analisis_en_Lean.C2_El_calculo_de_pi

variable {a : ℕ → ℝ}

-- 1ª demostración
-- ===============

example
  (ha : LimSuc a 5)
  : ∃ N, ∀ n ≥ N, 0 < a n :=
by
  obtain ⟨N, hN⟩ := ha 5 (by grind)
  -- N : ℕ
  -- hN : ∀ n ≥ N, |a n - 5| < 5
  refine ⟨N, by grind⟩

-- 2ª demostración
-- ===============

example
  (ha : LimSuc a 5)
  : ∃ N, ∀ n ≥ N, 0 < a n :=
by
  obtain ⟨N, hN⟩ := ha 5 (by norm_num)
  -- N : ℕ
  -- hN : ∀ n ≥ N, |a n - 5| < 5
  use N
  -- N : ℕ
  -- hN : ∀ n ≥ N, |a n - 5| < 5
  -- ⊢ ∀ n ≥ N, 0 < a n
  grind

-- 3ª demostración
-- ===============

example
  (ha : LimSuc a 5)
  : ∃ N, ∀ n ≥ N, 0 < a n :=
by
  obtain ⟨N, hN⟩ := ha 5 (by norm_num)
  -- N : ℕ
  -- hN : ∀ n ≥ N, |a n - 5| < 5
  use N
  -- N : ℕ
  -- hN : ∀ n ≥ N, |a n - 5| < 5
  -- ⊢ ∀ n ≥ N, 0 < a n
  intro n hn
  -- n : ℕ
  -- hn : n ≥ N
  -- ⊢ 0 < a n
  linarith [(abs_lt.mp (hN n hn)).1]

-- 4ª demostración
-- ===============

example
  (ha : LimSuc a 5)
  : ∃ N, ∀ n ≥ N, 0 < a n :=
by
  obtain ⟨N, hN⟩ := ha 5 (by norm_num)
  -- N : ℕ
  -- hN : ∀ n ≥ N, |a n - 5| < 5
  use N
  -- N : ℕ
  -- hN : ∀ n ≥ N, |a n - 5| < 5
  -- ⊢ ∀ n ≥ N, 0 < a n
  intro n hn
  -- n : ℕ
  -- hn : n ≥ N
  -- ⊢ 0 < a n
  have h1 := hN n hn
  -- h1 : |a n - 5| < 5
  have h2 := abs_lt.mp h1
  -- -5 < a n - 5 ∧ a n - 5 < 5
  have h3 := h2.1
  -- h3 : -5 < a n - 5
  exact (lt_add_iff_pos_left (-5)).mp h3

-- 5ª demostración
-- ===============

example
  (ha : LimSuc a 5)
  : ∃ N, ∀ n ≥ N, 0 < a n :=
by
  obtain ⟨N, hN⟩ := ha 5 (by norm_num)
  -- N : ℕ
  -- hN : ∀ n ≥ N, |a n - 5| < 5
  refine ⟨N, fun n hn => (lt_add_iff_pos_left (-5)).mp (abs_lt.mp (hN n hn)).1⟩

-- Aª demostración
-- ===============

example
  (ha : LimSuc a 5)
  : ∃ N, ∀ n ≥ N, 0 < a n :=
by
  have h0 : (0 : ℝ) < 5 := by norm_num
  specialize ha 5 h0
  -- ha : ∃ N, ∀ n ≥ N, |a n - 5| < 5
  choose N hN using ha
  -- N : ℕ
  -- hN : ∀ n ≥ N, |a n - 5| < 5
  use N
  -- ⊢ ∀ n ≥ N, 0 < a n
  intro n hn
  -- n : ℕ
  -- hn : n ≥ N
  -- ⊢ 0 < a n
  specialize hN n hn
  -- hN : |a n - 5| < 5
  rewrite [abs_lt] at hN
  -- hN : -5 < a n - 5 ∧ a n - 5 < 5
  linarith [hN.1]
\end{minted}

\begin{ejercicio}
  Demostrar que una sucesión cuyos términos superan 10 infinitas veces
  no puede converger a un límite de valor absoluto menor que 5; es
  decir, si la sucesión $a$ verifica que
     $$∀ N, ∃ n ≥ N, |aₙ| > 10$$
  entonces $a$ no puede tener un límite cuyo valor absoluto sea menor que
  5.
\end{ejercicio}

\begin{minted}{lean}
-- Demostración en lenguaje natural
-- ================================

-- Para demostrarlo por reducción al absurdo, supongamos que existe un
-- L ∈ ℝ tal que aₙ converge a L y
--    ∣L∣ < 5.                                                      (1)
-- Por la definición de límite, existe N ∈ ℕ tal que
--    ∀ n ≥ N, ∣aₙ − L∣ < 5                                         (2)
-- Por hipótesis sobre la sucesión, para ese mismo N existe tal que se
-- cumplen las siguientes dor relaciones:
--    n ≥ N                                                         (3)
--    ∣aₙ∣ > 10                                                     (4)
-- De (2) y (3) se tiene que
--    ∣aₙ − L∣ < 5                                                  (5)
-- Para obtener una contradicción basta probar que 0 < 0, como se
-- muestra a continuación:
--    0 < ∣aₙ∣              [por (4)]
--      = ∣(aₙ − L) + L∣
--      ≤ ∣aₙ − L∣ + ∣L∣
--      < 5 + 5             [por (5) y (1)]
--      = 10.

-- Demostraciones con Lean4
-- ========================

import Mathlib.Data.Real.Basic
import Mathlib.Tactic
import Analisis_en_Lean.C2_El_calculo_de_pi

variable {a : ℕ → ℝ}
variable {L : ℝ}

-- 1ª demostración
-- ===============

example
  (ha : ∀ N, ∃ n ≥ N, |a n| > 10)
  : ¬ ∃ L, |L| < 5 ∧ LimSuc a L :=
by
  intro ⟨L, hL1, hL2⟩
  -- L : ℝ
  -- hL1 : |L| < 5
  -- hL2 : LimSuc a L
  -- ⊢ False
  obtain ⟨N, hN⟩ := hL2 5 (by norm_num)
  -- N : ℕ
  -- hN : ∀ n ≥ N, |a n - L| < 5
  obtain ⟨n, hn1, hn2⟩ := ha N
  -- n : ℕ
  -- hn1 : n ≥ N
  -- hn2 : |a n| > 10
  apply lt_irrefl (10 : ℝ)
  -- ⊢ 10 < 10
  have h : |a n - L| < 5 := hN n hn1
  calc 10 < |a n|           := by exact hn2
        _ = |(a n - L) + L| := by rw [sub_add_cancel]
        _ ≤ |a n - L| + |L| := abs_add_le (a n - L) L
        _ < 5 + 5           := add_lt_add h hL1
        _ = 10              := by norm_num

-- 2ª demostración
-- ===============

example
  (ha : ∀ N, ∃ n ≥ N, |a n| > 10)
  : ¬ ∃ L, |L| < 5 ∧ LimSuc a L :=
by
  intro ⟨L, hL, hlim⟩
  -- L : ℝ
  -- hL : |L| < 5
  -- hlim : LimSuc a L
  -- ⊢ False
  obtain ⟨N, hN⟩ := hlim 5 (by norm_num)
  -- N : ℕ
  -- hN : ∀ n ≥ N, |a n - L| < 5
  obtain ⟨n, hn, hn10⟩ := ha N
  -- n : ℕ
  -- hn : n ≥ N
  -- hn10 : |a n| > 10
  have h1 : |a n - L| < 5 := hN n hn
  have h2 : |(a n - L) + L| ≤ |a n - L| + |L| := abs_add_le (a n - L) L
  rw [sub_add_cancel] at h2
  -- h2 : |a n| ≤ |a n - L| + |L|
  linarith

-- 3ª demostración
-- ===============

example
  (ha : ∀ N, ∃ n ≥ N, |a n| > 10)
  : ¬ ∃ L, |L| < 5 ∧ LimSuc a L :=
by
  intro ⟨L, hL1, hL2⟩
  -- L : ℝ
  -- hL1 : |L| < 5
  -- hL2 : LimSuc a L
  -- ⊢ False
  obtain ⟨N, hN⟩ := hL2 5 (by norm_num)
  -- N : ℕ
  -- hN : ∀ n ≥ N, |a n - L| < 5
  obtain ⟨n, hn1, hn2⟩ := ha N
  -- n : ℕ
  -- hn1 : n ≥ N
  -- hn2 : |a n| > 10
  apply lt_irrefl (10 : ℝ)
  -- ⊢ 10 < 10
  have h : |a n - L| < 5 := hN n hn1
  calc 10 < |a n|           := by exact hn2
        _ = |(a n - L) + L| := by ring_nf
        _ ≤ |a n - L| + |L| := abs_add_le _ _
        _ < 5 + 5           := by linarith
        _ = 10              := by norm_num

-- 4ª demostración
-- ===============

example
  (ha : ∀ N, ∃ n ≥ N, |a n| > 10)
  : ¬ ∃ L, |L| < 5 ∧ LimSuc a L :=
by
  intro ⟨L, hL1, hL2⟩
  -- L : ℝ
  -- hL1 : |L| < 5
  -- hL2 : LimSuc a L
  -- ⊢ False
  obtain ⟨N, hN⟩ := hL2 5 (by norm_num)
  -- N : ℕ
  -- hN : ∀ n ≥ N, |a n - L| < 5
  obtain ⟨n, hn1, hn2⟩ := ha N
  -- n : ℕ
  -- hn1 : n ≥ N
  -- hn2 : |a n| > 10
  apply lt_irrefl (10 : ℝ)
  -- ⊢ 10 < 10
  calc 10 < |a n|           := by exact hn2
        _ = |(a n - L) + L| := by grind
        _ ≤ |a n - L| + |L| := by grind
        _ < 5 + 5           := by grind
        _ = 10              := by grind

-- 5ª demostración
-- ===============

example
  (ha : ∀ N, ∃ n ≥ N, |a n| > 10)
  : ¬ ∃ L, |L| < 5 ∧ LimSuc a L :=
by
  intro ⟨L, hL1, hL2⟩
  -- L : ℝ
  -- hL1 : |L| < 5
  -- hL2 : LimSuc a L
  -- ⊢ False
  obtain ⟨N, hN⟩ := hL2 5 (by norm_num)
  -- N : ℕ
  -- hN : ∀ n ≥ N, |a n - L| < 5
  obtain ⟨n, hn, hn10⟩ := ha N
  -- n : ℕ
  -- hn : n ≥ N
  -- hn10 : |a n| > 10
  grind

-- 6ª demostración
-- ===============

example
  (ha : ∀ N, ∃ n ≥ N, |a n| > 10)
  : ¬ ∃ L, |L| < 5 ∧ LimSuc a L :=
by
  intro h
  -- h : ∃ L, |L| < 5 ∧ LimSuc a L
  -- ⊢ False
  obtain ⟨L, hL⟩ := h
  -- L : ℝ
  -- hL : |L| < 5 ∧ LimSuc a L
  have h1 : |L| < 5 := hL.1
  have ha' : LimSuc a L := hL.2
  have h0 : (0 : ℝ) < 5 := by norm_num
  specialize ha' 5 h0
  -- ha' : ∃ N, ∀ n ≥ N, |a n - L| < 5
  obtain ⟨N, hN⟩ := ha'
  specialize ha N
  -- ha : ∃ n ≥ N, |a n| > 10
  obtain ⟨n, hn⟩ := ha
  -- n : ℕ
  -- hn : n ≥ N ∧ |a n| > 10
  specialize hN n hn.1
  -- hN : |a n - L| < 5
  have ha : |a n| > 10 := hn.2
  have h2 : |a n| = |(a n - L) + L| := by ring_nf
  have h3 : |(a n - L) + L| ≤ |(a n - L)| + |L| := abs_add_le (a n - L) L
  linarith [h1, ha, h2, h3, hN]

-- Lemas usados
-- ============

variable (a b c d: ℝ)
#check (abs_add_le a b : |a + b| ≤ |a| + |b|)
#check (add_lt_add : a < b → c < d → a + c < b + d)
#check (lt_irrefl a : ¬a < a)
#check (sub_add_cancel a b : (a - b) + b = a)
\end{minted}

\begin{ejercicio}
Demostrar que si las sucesiones \(a\), \(c\) y \(e\) convergen a \(L\) y, para todo
\(n\),
   \[aₙ ≤ bₙ ≤ cₙ ≤ dₙ ≤ eₙ\]
enconces \(b\) y \(d\) también converge a \(L\).
\end{ejercicio}

\begin{minted}{lean}
import Mathlib.Data.Real.Basic
import Mathlib.Tactic
import Analisis_en_Lean.C6g_Teorema_del_emparedado

variable (a b c d e : ℕ → ℝ)
variable (L : ℝ)

-- 1ª demostración
-- ===============

example
  (ha : LimSuc a L)
  (hc : LimSuc c L)
  (he : LimSuc e L)
  (hab : ∀ n, a n ≤ b n)
  (hbc : ∀ n, b n ≤ c n)
  (hcd : ∀ n, c n ≤ d n)
  (hde : ∀ n, d n ≤ e n)
  : LimSuc b L ∧ LimSuc d L :=
by
  constructor
  · -- ⊢ LimSuc b L
    exact teorema_emparedado a b c L ha hc hab hbc
  · -- ⊢ LimSuc d L
    exact teorema_emparedado c d e L hc he hcd hde

-- 2ª demostración
-- ===============

example
  (ha : LimSuc a L)
  (hc : LimSuc c L)
  (he : LimSuc e L)
  (hab : ∀ n, a n ≤ b n)
  (hbc : ∀ n, b n ≤ c n)
  (hcd : ∀ n, c n ≤ d n)
  (hde : ∀ n, d n ≤ e n)
  : LimSuc b L ∧ LimSuc d L :=
⟨teorema_emparedado a b c L ha hc hab hbc,
 teorema_emparedado c d e L hc he hcd hde⟩
\end{minted}


%%% Local Variables:
%%% mode: LaTeX
%%% TeX-master: "Analisis_en_Lean"
%%% End:
